\chapter{Experimental Results and Discussion}

\section{Experimental Setup}

This section presents the experimental evaluation of the proposed
Temporal-State-Aware Routing Framework. The simulation experiments are
designed to evaluate routing reliability, recovery capability, robustness,
and scalability in dynamic quantum networks.

The simulation environment considers a quantum network consisting of
entanglement links with dynamically changing temporal states. Each link
is characterized by fidelity degradation, coherence budget, and expiration
state. The proposed framework is evaluated under different experimental
conditions including baseline comparison, ablation analysis, multi-seed
validation, regeneration budget analysis, and large-scale network testing.

The main evaluation metrics are:

\begin{itemize}
    \item Request acceptance ratio;
    \item Routing reliability;
    \item Link recovery capability;
    \item Stability across different random seeds;
    \item Scalability with increasing network size.
\end{itemize}


\section{Comparison with Baseline Routing Method}

To evaluate the effectiveness of the proposed framework, the proposed
Temporal-State-Aware Routing method is compared with a conventional
baseline routing strategy.

The comparison results demonstrate that considering temporal link quality
provides better routing decisions than static routing approaches. By
monitoring link aging and fidelity degradation, the proposed method avoids
selecting unreliable entanglement paths.

\begin{figure}[htbp]
\centering
\includegraphics[width=0.75\textwidth]
{figures/baseline_comparison.png}
\caption{Comparison between the proposed routing method and baseline approach.}
\label{fig:baseline}
\end{figure}

As shown in Figure~\ref{fig:baseline}, the proposed method achieves higher
routing reliability because routing decisions are continuously adapted
according to the current temporal state of entanglement links.


\section{Ablation Study}

An ablation study is performed to analyze the contribution of different
components in the proposed framework.

Three configurations are evaluated:

\begin{enumerate}
    \item Static baseline without temporal adaptation;
    \item Adaptive routing without regeneration;
    \item Complete proposed framework with temporal awareness and regeneration.
\end{enumerate}


\begin{figure}[htbp]
\centering
\includegraphics[width=0.75\textwidth]
{figures/ablation_comparison.png}
\caption{Ablation study showing the contribution of each framework component.}
\label{fig:ablation}
\end{figure}


The results indicate that temporal awareness improves route selection,
while regeneration provides additional recovery capability for expired
entanglement links.


\section{Multi-Seed Validation}

To verify the robustness of the proposed framework, experiments are
repeated using multiple random seeds.

\begin{figure}[htbp]
\centering
\includegraphics[width=0.75\textwidth]
{figures/multi_seed_comparison.png}
\caption{Performance stability under different random seeds.}
\label{fig:multiseed}
\end{figure}


The results show that the proposed framework maintains consistent
performance across different network realizations, demonstrating its
robustness and reliability.


\section{Regeneration Capability Analysis}

The impact of entanglement regeneration is investigated by changing the
regeneration budget.

\begin{figure}[htbp]
\centering
\includegraphics[width=0.75\textwidth]
{figures/regeneration_by_slot.png}
\caption{Temporal evolution of entanglement regeneration events.}
\label{fig:regeneration}
\end{figure}


The regeneration mechanism improves routing availability by restoring
expired links and increasing the probability of successful entanglement
distribution.


\section{Scalability Evaluation}

The scalability of the proposed framework is evaluated by increasing the
network size.

The results demonstrate that the proposed routing strategy maintains
stable performance when the network scale increases.

\begin{figure}[htbp]
\centering
\includegraphics[width=0.75\textwidth]
{figures/network_size_acceptance_ratio.png}
\caption{Scalability evaluation with increasing network size.}
\label{fig:scalability}
\end{figure}


\section{Summary}

The experimental results confirm that the proposed Temporal-State-Aware
Routing Framework improves quantum network reliability by jointly
considering temporal link conditions and regeneration capability.

The proposed approach achieves better routing performance, higher
robustness, and improved scalability compared with static routing
strategies.